\begin{theorem}
For each $\iota>0$, there is $D_0$ such that if $D\ge D_0$ and
$\mathcal H$ is a $3$-uniform $3$-simple hypergraph with maximum degree at
most $D$, then
\[
\chi'(\mathcal H)\le
\left(1-\frac19+\frac{1}{3^5}+\iota\right)3D.
\]
In particular, for sufficiently large $D$, every $3$-uniform $3$-simple
hypergraph with maximum degree at most $D$ has chromatic index at most
$2.6791D$.
\end{theorem}
\begin{proof}
Fix $\iota>0$.  Choose a positive real number $\alpha$ such that
$\alpha\le\iota$ and
\[
1-\frac19+\frac1{3^5}+\frac{\alpha}{2}<1;
\]
this is possible because $1-\frac19+\frac1{3^5}<1$.  Set
\[
\varepsilon=\alpha/2,\qquad
\theta=1-\frac19+\frac1{3^5}+\alpha/2 .
\]
Then $\varepsilon>0$ and $0<\theta<1$.

Apply \noderef{GeneralColoringTheorem} with $k=3$, this $\varepsilon$, and
this $\theta$.  It gives a threshold $D_{\mathrm{GCT}}$ such that every
$3$-uniform hypergraph $\mathcal H$ of maximum degree at most $D\ge
D_{\mathrm{GCT}}$ has chromatic index at most $(\theta+\varepsilon)3D$,
provided that for every $D'\ge\varepsilon D/3$, every sub-hypergraph
$H\subseteq\mathcal H$ of maximum degree at most $D'$, and every edge
$e\in E(H)$, one has
\[
b_{L(H),3D'}(e)\le\theta .
\]

Apply \noderef{ThreeUniformThreeSimpleBParameterEstimate} with
$\eta=\alpha/2$.  It gives a threshold $D_{\mathrm b}$ such that every
$3$-uniform $3$-simple hypergraph of maximum degree at most a natural
number $n\ge D_{\mathrm b}$ satisfies
\[
b_{L(H),3n}(e)\le
1-\frac19+\frac1{3^5}+\alpha/2=\theta
\]
for every edge $e$.

Choose $D_0$ large enough that $D_0\ge D_{\mathrm{GCT}}$ and
$\varepsilon D/3\ge D_{\mathrm b}+1$ for every $D\ge D_0$.  Now let
$D\ge D_0$, and let $\mathcal H$ be $3$-uniform, $3$-simple, and of maximum
degree at most $D$.  Fix any real $r=D'\ge\varepsilon D/3$ and any
sub-hypergraph $H\subseteq\mathcal H$ whose vertex degrees are at most $r$.
By \noderef{SubhypergraphInheritance}, $H$ is $3$-uniform and $3$-simple.
Put $n=\lfloor r\rfloor_+$.  Since $r\ge D_{\mathrm b}+1$, we have
$r\ge1$ and $n\ge D_{\mathrm b}$.
By \noderef{RealDegreeFloorBound}, $H$ has maximum degree at most $n$.
Apply the integer estimate above at $n$, and then apply
\noderef{LocalBParameterRealFloorTransfer} with $k=3$ to obtain, for every edge,
\[
b^{\mathbb R}_{L(H),3r}(e)\le b_{L(H),3n}(e)\le\theta.
\]
This verifies the real-scale local hypothesis of
\noderef{GeneralColoringTheorem}, applied with main scale $(D:\mathbb R)$.
Therefore
\[
\chi'(\mathcal H)\le(\theta+\varepsilon)3D
=\left(1-\frac19+\frac1{3^5}+\alpha\right)3D
\le
\left(1-\frac19+\frac1{3^5}+\iota\right)3D.
\]
This proves the asymptotic chromatic-index bound.

For the decimal conclusion,
\[
3\left(1-\frac19+\frac1{3^5}\right)
=3\cdot\frac89+\frac{3}{243}
=\frac83+\frac1{81}
=2.679012\ldots <2.6791.
\]
Choose a positive $\iota$ so small that
$3(1-\frac19+\frac1{3^5}+\iota)\le2.6791$, and use the first part with this
$\iota$.  For all sufficiently large $D$ every $3$-uniform $3$-simple
hypergraph of maximum degree at most $D$ then satisfies
\[
\chi'(\mathcal H)\le
\left(1-\frac19+\frac1{3^5}+\iota\right)3D
\le2.6791D,
\]
which is the stated decimal bound.
\end{proof}
